\documentclass[11pt]{article}
\usepackage[a4paper, top=18mm, bottom=18mm, left=20mm, right=20mm]{geometry}
\usepackage[T2A]{fontenc}
\usepackage[utf8]{inputenc}
\usepackage[ukrainian]{babel}
\usepackage{graphicx}
\usepackage[export]{adjustbox}
\usepackage{amsmath}
\usepackage{amssymb}
\graphicspath{{/app/Quiz/Scripts/2023/ProgAll/15BinTree/}{/app/Quiz/Scripts/2020/Mod2020_4/BinTree/}{/app/Quiz/Scripts/2021/Mod2021_3/BinTree/}}
\setlength{\parindent}{0pt}
\setlength{\parskip}{4pt}
\pagestyle{empty}
\begin{document}
%begin
{\large\textbf{Контрольна робота}}

\textbf{Варіант \No\,1}\hfill \textit{ПІБ:}\,\underline{\hspace{60mm}}
\vspace{4mm}

%beginitem
1. 
Скільки 2017-елементних підмножин множини натуральних від 1 до 20172017 містять рівно одне число, що більше 172016
%enditem

%beginitem
2. 

a,b,c=8,7,9
def f(a,b=6,c):
    print(a,b,c,end="")

f(a=4,1,a=3)
print(a,b,c)


%enditem

%beginitem
3. 

\begin{center}
	\adjustimage{max size={0.8\linewidth}{0.8\paperheight}}
	{../15BinTree/trees_exp/16.png}
\end{center}
Указати послідовність міток вершин дерева, зображеного на рис., що утвориться в результаті обходу дерева в прямому порядку. Мітки записувати через пробіл.\par Алгоритм починає роботу з кореня (на рис. це верхня вершина).
%answer:16 прямому

%enditem

%beginitem
4. %goodpath:Scripts/2018/Mod2018_1/01-good.txt
Відмітити все, що є однією правильною лексемою:
( 1 ) 3a

( 2 ) "1e3"

( 3 ) <+

( 4 ) {

%enditem

%beginitem
5. 
Записати ВИРАЗ, значенням якого є множина, що складається з квадратівцілих чисел від 16 до 2017 (включно), що не діляться на 5.
%enditem

%beginitem
6. Що буде виведено за виконання фрагменту коду?

%code
#include <iostream>

int f(int a, int b){
    int c = 26;
    if (b == 1)
        return 4;
    if (a <= 5)
         c = 6;
    else 
        c = 3;
    return c;
}

int main(){
    std::cout << f(8, 8);
    return 0;
}
%code

%enditem

%beginitem
7. Що буде виведено за виконання фрагменту коду?

%code
#include <iostream>
using namespace std;

bool f(int n){
    cout<<"f";
    return n;
}

int main(){
    cout<<(f(8) && f(2));
    return 0;
}
%code

%enditem

%beginitem
8. Що буде виведено за виконання фрагменту коду?

%code
print('R', end=' ')
m=11
while m>4:
    m+=-3
print(m)
%code


%enditem

%beginitem
9. Що буде виведено за виконання фрагменту коду?

%code
cout << 10 * 10 * 10 * 10;
%code

%enditem

%beginitem
10. 
try:
    res = 9/9//9*9
    if isinstance(res, bool):
        print(f'bool;{res}')
    elif isinstance(res, int):
        print(f'int;{res}')
    elif isinstance(res, float):
        print(f'float;{res:.2f}')
    else:
        print('???')
except:
    print('Z')

%enditem

%beginitem
11. 
a = 4
b = 6
c = 3

def f(a):
global c    a = 4
    b = 3
    c = 2
    return a + b + c

a = 7
b = 8
c = 5

try:
    print(f(b), a, b, c, sep=';')
except:
    print('ERR', a, b, c, sep=';')

%enditem

%beginitem
12. Що буде виведено за виконання фрагменту коду?

%code
print("7.0" == 8.0)
%code

%enditem

%beginitem
13. 
code
_c = 0

class C:
    _c = 1
    
    def __init__(self):
        self._c = 2
        _c = 3
        C._c = 4

obj = C()
try:
    print(_c, end=' ')
    print(obj._c, end=' ')
    print(C._c, end=' ')
    print(obj._c, end=' ')
    print(obj._C__c, end=' ')
except:
    print('ERR')
code

%enditem

%beginitem
14. Що буде виведено за виконання фрагменту коду?

%code
e = ('0','1','2','3','4','5','6','7','8','9')
print(e[2:2:2])
%code


%enditem

%beginitem
15. 
Рядки журналу через двокрапку містять ім'я користувача (ідентифікатор), час події,тип події, код події, наприклад
\begin{verbatim}
s="username : 2022-05-05 13:26 : ERROR : 404"
\end{verbatim}
у коді 
\begin{verbatim}
res = re.fullmatch('???', s) 
s = ''
code_str = ??? 
\end{verbatim}
чим треба замінити кожен з ???, щоб значенням code\_str був код події? 



%enditem

%beginitem
16. Що буде виведено за виконання фрагменту коду?

%code
print('R', end=' ')
e=('a','b','c','d','e',0,1,2,4)
print(e[11])
%code


%enditem

%beginitem
17. Що буде виведено за виконання фрагменту коду?

%code
a,b,c=4,6,3
def f(a):
global c    a=4
    b=3
    c=2
    return a+b+c

a,b,c=7,8,5
print(f(b),a,b,c)
%code

%enditem

%beginitem
18. 
e=('0','1','2','3','4','5','6','7','8','9'); print(e[2:2:2])
%enditem

%beginitem
19. %goodpath:Scripts/2019/Mod2019_1/Lex/03-good.txt
Відмітити все, що є коректним оператором С++:
( 1 ) 'a'a'

( 2 ) x=float("nan")

( 3 ) 'a'=3

( 4 ) F"1={x:5.2f};"

%enditem

%beginitem
20. 
x,y,z=4,1,3;  y,y,x,x=x,8,z,7;  print(x,y,z)
%enditem

%beginitem
21. 
def f(n): print("f", end=""); return n
   print(f(8) and f(2))
%enditem

%beginitem
22. 
m=4
   while m<11: m+=1
   print(m)
%enditem

%beginitem
23. 
code
c = 0
n = 1

class C:
    c = 2
    
    def __init__(self):
global c        self.c = 3
        c = 4
        C.c = 5

obj = C()
print(c, n, C.c, obj.c, sep=' ')
code

%enditem

%beginitem
24. Що буде виведено за виконання фрагменту коду?

%code
print('R', end=' ')
e = ('a','b','c','d','e',0,1,2,4)
print(e[11])
%code


%enditem

%beginitem
25. 
Значенням змінної \kw{a} є цілочисловий масив типу $numpy.ndarray$ розмірів $6{\times}6$. На скільки збільшиться сума елементів масиву \kw{a} після виконання
\begin{verbatim}
a[1] += 1
\end{verbatim}


Якщо код некоректний, то в якості відповіді зазначити ERROR.



%enditem

%beginitem
26. 
$a$ є цілочисловим масивом типу $numpy.ndarray$ розмірів $4{\times}4$. Сума елементів масиву $a$ дорівнює 37. Чому дорівнює сума елементів масиву $a$ після виконання 
\begin{verbatim}
a += 37
\end{verbatim}


Якщо код некоректний, то в якості відповіді зазначити ERROR.



%enditem

%beginitem
27. 
not6==6.0 and not8.0<=8 and 5>=True
%enditem

%beginitem
28. 
e=[10,11,12,13,14]; e[:7]=[1,2]; print(e)
%enditem

%beginitem
29. 
Змінні \kw{next} та \kw{prev} вузла списку позначають наступний та попередній за ним вузол відповідно. Починаючи з голови, у список записано послідовні цілі числа від~$-54$ до~$31$. Нехай змінна \kw{p} позначає вузол, що зберігає значення~$8$.
Яке число зберігається у вузлі \kw{p.prev.prev.prev.prev.prev} ?

%enditem

%beginitem
30. 
def f(arg, n):
    n = 4
    tmp = arg
    tmp[:7] = 'xy'
    return len(tmp)>3

try:
    e = [10,11,12,13,14]
    n = 6
    f(e, n)
    print(n, end=';')
    for el in e:
        print(el, end=';')
except:
    print('Z')

%enditem

%beginitem
31. 
e=[10,11,12,13,14]; e+=[2,3]; print(e)
%enditem

%beginitem
32. Що буде виведено за виконання фрагменту коду?
%code

x,y,z=4,1,3
y,y,x,x=x,8,z,7
print(x,y,z)
%code

%enditem

%beginitem
33. Що буде виведено за виконання фрагменту коду?

%code
print(not6==6.0 and not8.0<=8 and 5>=True)
%code

%enditem

%beginitem
34. 

print('R', end=' ')
e=[10,11,12,13,14]
e[:7]=[1,2]
print(e)


%enditem

%beginitem
35. 
def f(arg, n):
    n = 4
    tmp = arg
    tmp[:7] = 'xy'
    return len(tmp)>3

try:
    e = [10,11,12,13,14]
    n = 6
    f(e, n)
    print(n, end=';')
    for el in e:
        print(el, end=';')
except:
    print('ERR')

%enditem

%beginitem
36. Що буде виведено за виконання фрагменту коду?

%code
print("7.0"==8.0)
%code

%enditem

%beginitem
37. Що буде виведено за виконання фрагменту коду?

%code
#include <iostream>
using namespace std;

int f(int a, int b){
    int c = 26;
    if (b == 1)
        return 4;
    if (a <= 5)
         c = 6;
    else 
        c = 3;
    return c;
}

int main(){
    cout << f(8, 8);
    return 0;
}
%code

%enditem

%beginitem
38. 

\begin{center}
	\adjustimage{max size={0.8\linewidth}{0.8\paperheight}}
	{../BinTree/trees_exp/16.png}
\end{center}
Указати послідовність міток вершин дерева, зображеного на рис., що утвориться в результаті обходу дерева в прямому порядку. Мітки записувати через пробіл.\par Обхід дерева починається з кореня (на рис. це верхня вершина).
%answer:16 прямому

%enditem

%beginitem
39. Що буде виведено за виконання фрагменту коду?
%code
#include <iostream>
using namespace std;

int f(int &x, int &y){
    x = 7;
    y+= 8;
    return x;
}

int main(){
    int a = 2, b = 1;
    a = f(a, a);
    cout << a << ":" << b <<':';
    {
        int a = 5, b = 3;
        cout << ((a>6) && ((b+=1) > 6)) << ':';
        cout << a << ":" << b << ':';
    }
    {
        inta = 9;
        cout << a << ':';
    }
    cout << a << ":" << b << endl;
    return 0;
}
%code


%enditem

%beginitem
40. 
a,b,c=4,6,3
    def f(a):
      global c      a=4
      b=3
      c=2
      return a+b+c
    a,b,c=7,8,5
    print(f(b),a,b,c)
%enditem

%beginitem
41. 
В умовах попередньої задачі було виконано p->prev=p->prev->prev;

Чому дорівнює значення p->prev->prev->prev->prev->n ?
%enditem

%beginitem
42. Що буде виведено за виконання фрагменту коду?

%code
e=[10,11,12,13]
e+=[2,3]
print(e)
%code


%enditem

%beginitem
43. Що буде виведено за виконання фрагменту коду?

%code
9/9//9*9
%code

%enditem

%beginitem
44. 

print('R', end=' ')
m=11
while m>4:
    m+=-3
print(m)


%enditem

%beginitem
45. Що буде виведено за виконання фрагменту коду?

%code
print('R', end=' ')
e=[10,11,12,13,14]
e[:7]=[1,2]
print(e)
%code


%enditem

%beginitem
46. Що буде виведено за виконання фрагменту коду?

%code
a = 4
b = 6
c = 3

def f(a):
global c    a=4
    b+=3
    c=2
    return a+b+c

a, b, c = 7, 8, 5
try:
    print(f(b), a, b, c, sep=';')
except:
    print('Z', a, b, c, sep=';')
%code

%enditem

%beginitem
47. Що буде виведено за виконання фрагменту коду?

%code
#include <iostream>
using namespace std;

int a = 4, b = 6, c = 3;

int f(int &a){
int c;    a = 4;
    b -= 3;
    c = 2;
    return a + b + c;
}

int main(){
    inta = 7;
    intb = 8;
    intc = 5;
    cout << f(b) << ':';
    cout << a << ':' << b << ':' << c;
    return 0; 
}
%code

%enditem

%beginitem
48. 
a,b,c=4,6,3
    def f(a):
      global c      a=4
      b+=3
      c=2
      return a+b+c
    a,b,c=7,8,5
    print(f(b),a,b,c)
%enditem

%beginitem
49. Що буде виведено за виконання фрагменту коду?

%code
try:
    m = 11
    while m>4:
        m += -3
    print(m, end=';')
except:
    print('Z')

%code

%enditem

%beginitem
50. 

\begin{center}
	\adjustimage{max size={0.8\linewidth}{0.8\paperheight}}
	{../BinTree/trees_exp/16.png}
\end{center}
Указати послідовність міток вершин дерева, зображеного на рис., що утвориться в результаті обходу дерева в прямому порядку. Мітки записувати через пробіл.\par Алгоритм починає роботу з кореня (на рис. це верхня вершина).
%answer:16 прямому

%enditem

%beginitem
51. 
try:
    m = 11
    while m>4:
        m += -3
    print(m, end=';')
except:
    print('Z')

%enditem

%beginitem
52. Що буде виведено за виконання фрагменту коду?

%code
def f(a,b):
    c=26
    if b==1:
        return 4
    if a<=5:
         c=6
    else: 
        c=3
    return c

print(f(8,8))
%code

%enditem

%beginitem
53. Що буде виведено за виконання фрагменту коду?

%code
#include <iostream>
using namespace std;

int f(int c){
    int y = 26;
    if (c == 1) 
        return 4;
    if (c <= 5)
         y = 6;
    else
         y = 3;
    return y;
}

int main(){
    cout << f(8);
    return 0;
}
%code


%enditem

%beginitem
54. %goodpath:Scripts/2018/Mod2018_1/03-good.txt
Відмітити все, що є коректним оператором С++:
( 1 ) cout<<(3<>5);

( 2 ) int x='a'*'c';

( 3 ) int n,n;

( 4 ) cout<<'\n';

%enditem

%beginitem
55. Що буде виведено за виконання фрагменту коду?

%code
def f(c):
    y=26
    if c==1: 
        return 4
    if c<=5:
         y=6
    else:
         y=3
    return y

print(f(8))
%code

%enditem

%beginitem
56. Що буде виведено за виконання фрагменту коду?

%code
cout << (!6==6.0 && !8.0<= 8 && 5>=true);
%code

%enditem

%beginitem
57. Що буде виведено за виконання фрагменту коду?

%code
print('R', end=' ')
e = ('0','1','2','3','4','5','6','7','8','9')
print(e[2:2:2])
%code


%enditem

%beginitem
58. 

print('R', end=' ')
e=[10,11,12,13,14]
e+=[2,3]
print(e)


%enditem

%beginitem
59. Що буде виведено за виконання фрагменту коду?

%code
print('R', end=' ')
e=[10,11,12,13,14]
e+=[2,3]
print(e)
%code


%enditem

%beginitem
60. Що буде виведено за виконання фрагменту коду?

%code
print(9/9//9*9)
%code

%enditem

%beginitem
61. 
try:
    m = 11
    while m>4:
        m += -3
    print(m, end=';')
except:
    print('ERR')

%enditem

%beginitem
62. 
try:
    for e in range(2, 2, 2):
        if e>7:
            continue
        print(e, end=';')
        if e>7:
            break
    else:
        print(e,end=';')
    print(e)
except:
    print('ERR')

%enditem

%beginitem
63. 
"7.0"==8.0
%enditem

%beginitem
64. Що буде виведено за виконання фрагменту коду?
%Оригинал был утерян. Требует отладки!
%code
_c = 0

class C:
    _c = 1
    
    def __init__(self):
        self._c = 2
        _c = 3
        C._c = 4

obj = C()
obj._c = 5
try:
    print(_c, end=' ')
    print(obj._c, end=' ')
    print(C._c, end=' ')
    print(obj._c, end=' ')
    print(obj._C__c)
except:
    print('ERR')
%code

%enditem

%beginitem
65. 
У папці X77 розташовано проект, до якого входить синаксично правильна одиниця трансляції 1.cpp з рядком\par
{\#}include "2.h"
\parФайла 2.h у папці X77 нема, але він є в папці Y2. Що треба зробити, щоб одиниця трансляції 1.cpp, могла успішно відкомпілюватися?Переміщувати, копіювати та змінювати файли з кодом заборонено.\par\vspace{1.5cm}
%enditem

%beginitem
66. 

print('R', end=' ')
for e in range(2,2,2):
    if e>7:
        continue
        print(e, end=' ')
    if e>7:
        break
    else:
        print(e, end=' ')
    print(e)

%enditem

%beginitem
67. Що буде виведено за виконання фрагменту коду?

%code
m=4
while m<11:
    m+=1
print(m, end=' ')
%code


%enditem

%beginitem
68. 
Записати ВИРАЗ, значенням якого є словник, ключами якого є цілі числа від 16 до 2017 (включно), що не діляться на 5, а ключам відповідають їх квадрати 
%enditem

%beginitem
69. 
Написати програму, що запитує у користувача значення дійсних параметрів $x$ та $y$, обчислює для них значення виразу 
$$\frac{\sqrt x + 22}{(\arctg x + 0.7) (y + 93)}$$ 
та повiдомляє введенi данi та результат обчислення.
 Оцінюється правильність та якість коду.


%enditem

%beginitem
70. %goodpath:Scripts/2019/Mod2019_1/Lex/01-good.txt
Відмітити все, що є однією правильною лексемою:
( 1 ) a++

( 2 ) """1"2"""

( 3 ) <+

( 4 ) {

%enditem

%beginitem
71. 
Змінні $next$ та $prev$ вузла списку позначають наступний та попередній за ним вузол відповідно. Починаючи з голови, у список записано послідовні цілі числа від $-54$ до $31$. Нехай змінна $p$ позначає вузол, що зберігає значення $8$.
Яке значення зберігається у вузлі $p.prev.prev.prev.prev.prev$ ?

%enditem

%beginitem
72. 
for e in range(2,2,2):
    if e>7:
        continue
        print(e,end=' ')
    if e>7:
        break
    else:
        print(e,end=' ')
    print(e)

%enditem

%beginitem
73. 
(Задача за частиною Пашка А.О. Наводиться в авторській редакції.)\par {\hspace {0.5 cm}}³дбудеться деяка подія $W$ чи ні залежить від двох факторів $A$ та $X$. За результатами статистичного експерименту (100 раз для кожного значення фактора $A$,  200 раз для кожного значення фактора $X$) подія $W$ відбулася:\par {\hspace {0.5 cm}}$(n+17)$ раз для значення $A(a)$, $(n+21)$ раз для значення $A(b)$, $(n+50)$ раз для значення $A(c)$, $(n+57)$ раз для значення $A(d)$, та відповідно $(2n+100)$ раз для значення $X(x)$, $(4n+17)$ раз для значення $X(y)$, $(n+150)$ раз для значення $X(z)$.\par
{\hspace {0.5 cm}}Знайти ймовірність того що подія $W$ не відбулася (умовна ймовірність) при значеннях факторів $A(b)$, $X(y)$. Фактори є незалежними, $n$ – номер цього екзаменаційного білета.\par
%enditem

%beginitem
74. 
for e in range(2,2,2):
    if e>7:
        continue
        print(e,end=' ')
    if e>7:
        break
else:
    print(e, end=' ')
print(e, end=' ')

%enditem

%beginitem
75. 
Записати ВИРАЗ, значенням якого є список, що складається з квадратівцілих чисел від 16 до 2017 (включно), що не діляться на 5, записаних в оберненому порядку.
%enditem

%beginitem
76. 

def f(c):
    y=26
    if c==1: 
        return 4
    if c<=5:
         y=6
    else:
         y=3
    return y

print(f(8))

%enditem

%beginitem
77. 
Записати ВИРАЗ, значенням якого є словник, ключами якого є цілі числа від 16 до 2017 (включно), що не діляться на 5, а ключам відповідають їх квадрати 
%enditem

%beginitem
78. 
m=11
   while m>4: m+=-3
   print(m)
%enditem

%beginitem
79. Що буде виведено за виконання фрагменту коду?
%code
for e in range(2,2,2):
    if e>7:
        continue
        print(e,end=' ')
    if e>7:
        break
else:
    print(e, end=' ')
print(e, end=' ')
%code

%enditem

%beginitem
80. 
Проект складається з од���ї одиниці трансляції 1.cpp, яка починається таким фрагментом\par
long f(){ return myFactorial();}\par
Функція myFactorial(); означена в 2.cpp.\par
Що треба зробити для того, щоб 1.cpp успішно відкомпілювався?\par\vspace{1.5cm}
Що треба зробити для того, щоб за проектом зібрався виконуваний код?\par

%enditem

%beginitem
81. Що буде виведено за виконання фрагменту коду?

%code
cout << (!6 == 6.0 && !8.0 <= 8 && 5 >= true);
%code

%enditem

%beginitem
82. 
Записати ВИРАЗ, значенням якого є список, що складається з квадратівцілих чисел від 16 до 2017 (включно), що не діляться на 5, записаних в оберненому порядку.
%enditem

%beginitem
83. Що буде виведено за виконання фрагменту коду?

%code
for e in range(7,7+5,2):
    if e>7:
        break
    print(e, end=' ')
    e=7
    if e>7:
        break
else:
    print(e, end=' ')
print(e, end=' ')
%code

%enditem

%beginitem
84. 
try:
    res = 1**1.0*True
    if isinstance(res, bool):
        print(f'bool;{res}')
    elif isinstance(res, int):
        print(f'int;{res}')
    elif isinstance(res, float):
        print(f'float;{res:.2f}')
    else:
        print('???')
except:
    print('ERR')

%enditem

%beginitem
85. Що буде виведено за виконання фрагменту коду?

%code
print(6is6.0is8.0is8)
%code

%enditem

%beginitem
86. Що буде виведено за виконання фрагменту коду?

%code
1**1.0*True
%code

%enditem

%beginitem
87. 
Знайти кількість відношень на n-елементній множині, які нерефлексивні та несиметричні
%enditem

%beginitem
88. Що буде виведено за виконання фрагменту коду?
%code
if (16 == 16)
    cout << "q";
else
    cout << "q";
%code


%enditem

%beginitem
89. Що буде виведено за виконання фрагменту коду?

%code
"7.0"==8.0
%code

%enditem

%beginitem
90. 

a,b,c=4,6,3
def f(a):
global c    a=4
    b+=3
    c=2
    return a+b+c

a,b,c=7,8,5
print(f(b),a,b,c)

%enditem

%beginitem
91. Що буде виведено за виконання фрагменту коду?

%code
cout << (6 == 6.0 <= 8.0 >= 8);
%code

%enditem

%beginitem
92. 
a = 4
b = 6
c = 3

def f(a):
global c    a = 4
    b = 3
    c = 2
    return a + b + c

a = 7
b = 8
c = 5

try:
    print(f(b), a, b, c, sep=';')
except:
    print('Z', a, b, c, sep=';')

%enditem

%beginitem
93. Що буде виведено за виконання фрагменту коду?

%code
print('R', end=' ')
m=4
while m<11:
    m+=1
print(m)
%code


%enditem

%beginitem
94. 
try:
    res = not6==6.0 and not8.0<=8 and 5>=True
    if isinstance(res, bool):
        print(f'bool;{res}')
    elif isinstance(res, int):
        print(f'int;{res}')
    elif isinstance(res, float):
        print(f'float;{res:.2f}')
    else:
        print('???')
except:
    print('Z')

%enditem

%beginitem
95. 
try:
    res = "7.0"==8.0j
    if isinstance(res, bool):
        print(f'bool;{res}')
    elif isinstance(res, int):
        print(f'int;{res}')
    elif isinstance(res, float):
        print(f'float;{res:.2f}')
    else:
        print('???')
except:
    print('ERR')

%enditem

%beginitem
96. 
Написати програму, що запитує у користувача значення дійсних параметрів $x$ та $y$, обчислює для них значення виразу 
$$\frac{\sqrt x + 22}{(\arctg x + 0.7) (y + 93)}$$ 
та повiдомляє введенi данi та результат обчислення.


%enditem

%beginitem
97. Що буде виведено за виконання фрагменту коду?

%code
cout << 9 * 9 * 9 * 9;
%code

%enditem

%beginitem
98. 
a = 4
b = 6
c = 3

def f(a):
global c    a = 4
    b += 3
    c = 2
    return a + b + c

a = 7
b = 8
c = 5

try:
    print(f(b), a, b, c, sep=';')
except:
    print('ERR', a, b, c, sep=';')

%enditem

%beginitem
99. Що буде виведено за виконання фрагменту коду?

%code
print('R', end=' ')
e=('0','1','2','3','4','5','6','7','8','9')
print(e[2:2:2])
%code


%enditem

%beginitem
100. 

try:
    print(4,end="")
    print(int("c6"),end="")
    print(3,end="")
except ZeroDivisionError: print(7,end="")
except BaseException: print(8,end="")
else: print(5,end="")
finally: print(9,end="")

%enditem

%end
\newpage
\end{document}

